\section{Methodology}
\label{sec:methodology}

% Subsections still to be written by their owners:
%   - Solver--Critic--Solver collection protocol and activation capture
%   - Sparse autoencoder training

\subsection{Behavioral labels}
\label{sec:method-labels}

Every episode is assigned one of the four behavior labels defined in
Section~\ref{sec:eda}: Adoption, Retention, Agreement, or Out. Only Adoption
and Retention enter the binary target used by the probes: Adoption is the
positive class and Retention the negative class. The labels come from two
sources, which we keep distinct throughout the paper.

\paragraph{Rule-based reference labels.} The labels reported in
Section~\ref{sec:eda} were produced by deterministic rules that compare A1,
the Critic's recommendation (C1), and A2 using normalized text, answer
aliases, and explicit handling of refusals, letter-only replies, citations,
and hybrid answers. These rules were audited against 46 team-reviewed
episodes and are frozen. We use them as the reference baseline and for every
analysis of the original MuSiQue collection.

\paragraph{LLM-derived training labels.} For the combined MuSiQue and
synthetic collection, the training targets are labels generated by Claude
Opus~5, accessed through Amazon Bedrock with a fixed prompt
(\texttt{claude\_binary\_v3\_candidates}), medium reasoning effort, and an
8,000-token output limit. For each episode the labeler receives the
question, A1, the Critic's recommended answer, and A2, presented in the same
candidate order the Solver saw. It returns one of the four labels with a
one-sentence reason. Episodes whose Critic gave no committed answer are
labeled Out by rule. API, parsing, or truncation failures are recorded as
technical exclusions and are never counted as a behavioral label. These are
model-generated labels, not independently verified ground truth.

\paragraph{Agreement with the reference labels.} Before adopting the
LLM-derived labels, we compared them with the reference labels on a
250-episode sample from the training partition of the original collection
(Table~\ref{tab:label-agreement}). On the Adoption--Retention decision that
defines the probe target, the two sources agree in every case (100 of 100
stratified episodes). Almost all disagreement falls in the reference
labels' Out class: the LLM labeler assigns a behavioral label to 47 of the
50 episodes that the rules leave out, many of them long-sentence or
letter-only replies that the rules exclude conservatively. On the 46
team-reviewed episodes, the LLM labels match the team's judgment in 38
cases, compared with 26 for the rule-based labels. The sample ran with no
technical failures, at about 2.3 seconds and \$0.006 per episode.

\begin{table}[htbp]
    \centering
    \caption{LLM-derived labels against the rule-based reference labels on a
    stratified 200-episode evaluation set (50 per reference label) from the
    training partition.}
    \label{tab:label-agreement}
    \begin{tabular}{lrrrr}
        \toprule
        Reference label & \multicolumn{4}{c}{LLM-derived label} \\
        \cmidrule(lr){2-5}
         & Adoption & Retention & Agreement & Out \\
        \midrule
        Adoption  & 50 &  0 &  0 & 0 \\
        Retention &  0 & 50 &  0 & 0 \\
        Agreement &  0 &  1 & 49 & 0 \\
        Out       & 23 & 11 & 13 & 3 \\
        \bottomrule
    \end{tabular}
\end{table}

Cohen's $\kappa$ on this set is 0.68, and accuracy re-weighted to the
training partition's label mix is 0.87. The lower agreement on Out reflects
a difference in how ambiguous replies are treated, not a disagreement about
which candidate the Solver chose.

\subsection{Aligning labels with activations}
\label{sec:method-alignment}

Labels are exported together with the frozen partition manifest as a
package that the probe and feature-selection code read directly. The export
refuses an episode labeled twice, an episode in a labeled partition that has
no label, and any label whose partition or activation index disagrees with
the collection. It also refuses a label outside the four classes. Technical
exclusions are listed with their reasons, so every collected episode is
either labeled or explicitly excluded. Partitions that are not labeled stay
in the manifest without labels and cannot be read by accident. On the
original collection, the export joins 22,130 episodes to their nine-layer
activations with every row index verified.

\subsection{Linear probes}
\label{sec:method-probes}

\paragraph{Inputs and target.} For episode $i$ and layer $\ell$, let
$\mathbf{x}_i^{(\ell)} \in \mathbf{R}^{2560}$ be the output of decoder layer
$\ell$ at the final token of the Solver's revision prompt, captured during the
prefill pass. At that point the Solver has read the Critic's recommendation but
has not yet written A2. The target is $y_i = 1$ for Adoption and $y_i = 0$ for
Retention. Agreement, Out, and technical exclusions are removed before fitting
and are never counted as Retention.

\paragraph{Probe.} Each layer gets its own logistic probe on standardized
activations, $\tilde{\mathbf{x}} = (\mathbf{x} - \mu)/\sigma$, where the mean
$\mu$ and standard deviation $\sigma$ of each unit come from the training rows.
The probe predicts
\begin{equation}
  p(y = 1 \mid \mathbf{x}) =
  \frac{1}{1 + \exp\{-(\mathbf{w}^{\top}\tilde{\mathbf{x}} + b)\}}
  \label{eq:probe}
\end{equation}
and is fit by minimizing the class-weighted, $L_2$-regularized log loss
\begin{equation}
  \frac{1}{2}\|\mathbf{w}\|^{2}
  + C \sum_{i} a_{y_i}\,
  \mathcal{L}\bigl(y_i,\ \mathbf{w}^{\top}\tilde{\mathbf{x}}_i + b\bigr),
  \qquad a_c = \frac{n}{2\,n_c},
  \label{eq:probe-loss}
\end{equation}
where $\mathcal{L}$ is the logistic loss and $n_c$ is the number of training
rows in class $c$, so both classes carry equal weight. The regularization
strength $C \in \{10^{-3}, 10^{-2}, \ldots, 10^{2}\}$ is chosen by three-fold
cross-validation inside the training partition, maximizing balanced accuracy,
with ties going to the smaller $C$. The probe is then refit on all training
rows. Folds are stratified by class and grouped by question family, so related
variants of a question never fall on both sides of a split.

\paragraph{Evaluation.} Layers are compared on the validation partition by
balanced accuracy, $\mathrm{BA} = (\mathrm{TPR} + \mathrm{TNR})/2$, and by the
AUROC of the score $s_i = \mathbf{w}^{\top}\tilde{\mathbf{x}}_i + b$, oriented
so that a larger score means Adoption. Each metric has a 95\% interval from
1{,}000 bootstrap resamples of whole question families. Every layer and
baseline uses the same resamples, so differences between them are paired. The
selected layer is the one with the highest validation balanced accuracy; if two
layers tie, none is selected. Validation results are also broken down by
dataset source, correctness cell (whether A1 and the Critic were right),
candidate order, and hop type. A subgroup is reported only when it contains
both classes.

\paragraph{Controls.} A probe can score well by reading something other than
the Solver's decision. Table~\ref{tab:probe-controls} lists the controls, each
fit and scored with the same procedure. The permuted-label control follows
Hewitt and Liang~\cite{hewitt2019control}: training labels are shuffled across
questions, keeping the class counts, and the probe is refit and scored on the
true validation labels. Over 100 permutations, a layer's selectivity is its
observed metric minus the mean of this null distribution.

\begin{table}[t]
  \centering
  \small
  \caption{Controls for the layer-wise probes. Each is fit on the training
  partition and scored on validation exactly like the activation probe.}
  \label{tab:probe-controls}
  \begin{tabular}{@{}p{0.2\linewidth}p{0.42\linewidth}p{0.32\linewidth}@{}}
    \toprule
    Control & Input & Rules out \\
    \midrule
    Visible text & TF-IDF (1--2-grams) of the question, A1, the Critic's
      recommended answer, and the candidate order: what the Solver sees
      & A decision readable from the prompt text alone \\
    Metadata & One-hot candidate order, correctness cell, dataset source,
      and hop type & Position, difficulty, or source effects \\
    Pre-feedback & The same probe on the layer's activation at the end of
      the first prompt, before A1 and before any feedback
      & Question difficulty rather than a response to feedback \\
    Permuted labels & Activations with the training labels shuffled
      (100 permutations) & A probe flexible enough to fit any labels
      (selectivity) \\
    \bottomrule
  \end{tabular}
\end{table}

\paragraph{Feature selection.} The same probe is fit on the latents of a sparse
autoencoder at the selected layer. With encoder $E$ and decoder $D$, the latent
vector is $\mathbf{z} = E(\mathbf{x}) \in \mathbf{R}^{m}$, and only latents that
fire on the training rows are used. Because the probe is linear in standardized
inputs, the SHAP value of latent $j$ for row $i$ is $w_j \tilde{z}_{ij}$
\cite{lundberg2017shap}. Latents are ranked on the training rows by their mean
absolute SHAP value,
\begin{equation}
  \bar{\phi}_j = |w_j| \cdot \frac{1}{n}\sum_{i=1}^{n} |\tilde{z}_{ij}|.
  \label{eq:shap}
\end{equation}
The targeted set $S$ is the ten highest-ranked latents with $w_j > 0$, the
latents whose activity pushes the probe toward Adoption. Ranking on training
rows keeps the choice of features independent of the rows on which their
effects are measured.

\paragraph{Raw-activation basis.} To test whether the autoencoder adds anything
beyond individual neurons~\cite{cunningham2023sparse,marks2025circuits}, the
same selection is run on the raw activations through an identity basis,
$E(\mathbf{x}) = \mathbf{x} - \mu$ and $D(\mathbf{z}) = \mathbf{z} + \mu$, with
$\mu$ the training mean. Setting a coordinate of $\mathbf{z}$ to zero then
replaces that neuron with its training mean (mean ablation), and the targeted
set is the ten neurons the probe ranks highest toward Adoption.
Figure~\ref{fig:alg-probes} summarizes the procedure.

\begin{figure}[t]
  \centering
  \fbox{\begin{minipage}{0.94\linewidth}
  \small
  \begin{enumerate}
    \item For each layer $\ell$: standardize $\mathbf{x}^{(\ell)}$ with training
      statistics, choose $C$ by family-grouped inner cross-validation on the
      training rows, refit, and score balanced accuracy and AUROC on validation
      with family-bootstrap intervals.
    \item Fit the controls in Table~\ref{tab:probe-controls} the same way, and
      compute each layer's selectivity against 100 label permutations.
    \item Select the layer $\ell^{*}$ with the highest validation balanced
      accuracy (none if two layers tie).
    \item For each basis at $\ell^{*}$ (each autoencoder and the raw-activation
      basis): encode the training rows, keep the live latents, fit the probe,
      rank the latents by Equation~\ref{eq:shap}, and take the top ten with
      $w_j > 0$ as $S$.
    \item Freeze $S$, the probe weights, and the matched random feature sets
      (Section~\ref{sec:method-causal}) before any intervention is run.
  \end{enumerate}
  \end{minipage}}
  \caption{Probe training and feature selection. Only the training and
  validation partitions are read; the test and intervention partitions stay
  sealed.}
  \label{fig:alg-probes}
\end{figure}

\subsection{Causal interventions}
\label{sec:method-causal}

A probe shows that information is present, not that the Solver uses it. To test
use, we replay episodes from the intervention partition, which no probe or
feature selection has read, and change the activation where it was measured.
For each eligible episode, the stored A1 and Critic feedback go back through
the Solver's revision prompt with the same model, revision, prompt, and greedy
decoding as in collection, and a new A2 is generated. The edit is applied once,
at layer $\ell$ and the final prompt token, during the prefill. Earlier
positions and later decoding steps are not modified. Before any edited
condition runs, the unedited replay is checked against collection: it must
reproduce the stored label for at least 98\% of episodes, or the run stops.

\paragraph{Edits.} Let $\mathbf{z} = E(\mathbf{x})$, and let
$\mathbf{z}_{S \leftarrow 0}$ be $\mathbf{z}$ with the latents in $S$ set to
zero. The contribution of $S$ to the reconstruction is
$\mathbf{c}_S = D(\mathbf{z}) - D(\mathbf{z}_{S \leftarrow 0})$.
Table~\ref{tab:edits} lists the edits. Ablation removes $\mathbf{c}_S$, but the
size of that change depends on how strongly the features in $S$ happen to fire.
In our preliminary runs, effects appeared only for large edits, and targeted
edits were no stronger than random edits of similar size. The norm-matched edit
therefore fixes the size relative to the activation and keeps only the
direction of $\mathbf{c}_S$:
\begin{equation}
  \Delta = r\,\|\mathbf{x}\|\,\frac{\mathbf{c}_S}{\|\mathbf{c}_S\|},
  \qquad r \in \{\pm 0.1, \pm 0.25, \pm 0.5, \pm 1\},
  \label{eq:norm-edit}
\end{equation}
where $r < 0$ suppresses what $S$ contributes and $r > 0$ amplifies it. Every
feature set, autoencoder, and episode then receives an edit of the same size.
When no latent in $S$ fires, $\mathbf{c}_S = 0$, and the edit is recorded as a
no-op and reported separately. For the other edits, which change $\mathbf{z}$ to
$\mathbf{z}'$, the edit is $\Delta = D(\mathbf{z}') - D(\mathbf{z})$. It is
applied either inside the reconstruction, $\mathbf{x}' = D(\mathbf{z}) + \Delta$,
or with the error-preserving splice, $\mathbf{x}' = \mathbf{x} + \Delta$, which
keeps the autoencoder's reconstruction error, so that the unedited control is
the original activation.

\begin{table}[t]
  \centering
  \small
  \caption{Interventions on the final-token activation. $\mathrm{ref}_j$ is the
  mean activation of latent $j$ where it fires on the training rows;
  $\mathbf{u}$ is a fixed random direction.}
  \label{tab:edits}
  \begin{tabular}{@{}p{0.2\linewidth}p{0.36\linewidth}p{0.38\linewidth}@{}}
    \toprule
    Edit & Change & Purpose \\
    \midrule
    Ablation & $z_j \leftarrow 0$ for $j \in S$ & Remove the targeted features \\
    Steering ($s$) & $z_j \leftarrow s\,\mathrm{ref}_j$ for $j \in S$
      & Set the features to a multiple of their typical active value \\
    Scaling ($p$) & $z_j \leftarrow (1 + p)\,z_j$ for $j \in S$, $p \ge -1$
      & Shrink or grow the features in proportion \\
    Norm-matched ($r$) & Equation~\ref{eq:norm-edit}
      & Same edit size for every feature set \\
    Random direction ($r$) & $\mathbf{x}' = \mathbf{x} + r\,\|\mathbf{x}\|\,
      \mathbf{u}/\|\mathbf{u}\|$ & Edit of the same size that does not use the
      autoencoder \\
    \bottomrule
  \end{tabular}
\end{table}

\paragraph{Controls.} Each targeted edit is paired with the same edit applied
to $K$ random feature sets of the same size. For each targeted latent, a random
latent is drawn from the non-targeted latents that fire within a factor of 1.25
as often on the training rows, among the 50 nearest in typical decoded
contribution, $\mathrm{RMS}(z_j)\,\|\mathbf{d}_j\|$, where $\mathbf{d}_j$ is
latent $j$'s decoder column. For the raw-activation basis, the decoder columns
are unit vectors, so this matches on each neuron's standard deviation. Three
fixed random directions add a control that does not depend on the autoencoder.
% K = 100 for the confirmatory run (Instructor Review #6); the earlier
% edit-size configurations used K = 5.

\paragraph{Estimands and tests.} Let $Y_i(c) = 1$ when the A2 regenerated under
condition $c$ is labeled Adoption by the rule-based labels
(Section~\ref{sec:method-labels}), and $0$ otherwise. For a targeted condition
$T$ and its unedited control, the effect is the paired change in the Adoption
rate,
\begin{equation}
  \delta_{T} = \frac{1}{n}\sum_{i=1}^{n} \bigl[Y_i(T) - Y_i(0)\bigr],
  \label{eq:delta}
\end{equation}
and the comparison that isolates the features is targeted against random,
\begin{equation}
  \delta_{T-R} = \frac{1}{n}\sum_{i=1}^{n}
  \Bigl[Y_i(T) - \frac{1}{K}\sum_{k=1}^{K} Y_i(R_k)\Bigr],
  \label{eq:targeted-vs-random}
\end{equation}
where $R_1, \ldots, R_K$ are the matched random sets at the same edit. Intervals
come from 1{,}000 bootstrap resamples of questions. For $\delta_{T-R}$, a second
interval also resamples the random sets, so it includes the uncertainty from
which sets were drawn. The $p$-value is a two-sided paired sign-flip test,
$p = (1 + \#\{b : |\delta^{(b)}| \ge |\delta|\})/(1 + B)$ with
$B = 10{,}000$, where each $\delta^{(b)}$ flips the sign of every question's
difference at random. Targeted-versus-random comparisons are Holm-adjusted
\cite{holm1979simple} within each edit family (ablation, steering, scaling,
norm-matched), with one test per experiment (autoencoder, feature subset, and
splice) and edit strength. Each row also
records the log-probability of the first generated token, a continuous measure
beside the binary label. Figure~\ref{fig:alg-causal} summarizes the procedure.
% #193: report Retention and third-answer rates per condition beside Adoption,
% and add the raw-neuron and probe-direction arms, once implemented.

\begin{figure}[t]
  \centering
  \fbox{\begin{minipage}{0.94\linewidth}
  \small
  \begin{enumerate}
    \item Load the frozen targeted set $S$, the random sets $R_1, \ldots, R_K$,
      and the autoencoder from feature selection. Nothing is fit on the
      intervention partition.
    \item Replay every episode without an edit; stop unless at least 98\% of
      the regenerated A2s keep their stored label.
    \item For each episode $i$, feature set ($S$, each $R_k$, each random
      direction), and edit strength: at layer $\ell$ and the final prompt
      token, replace $\mathbf{x}$ by $\mathbf{x}'$ (Table~\ref{tab:edits})
      during the prefill, regenerate A2 greedily, label it, and record
      $Y_i$, the first-token log-probability, and the applied edit size.
    \item Estimate $\delta_T$ and $\delta_{T-R}$ with bootstrap intervals and
      sign-flip $p$-values, and Holm-adjust within each edit family.
  \end{enumerate}
  \end{minipage}}
  \caption{Causal replay. The edit is applied only at the position where the
  activations, probes, and features were measured.}
  \label{fig:alg-causal}
\end{figure}
